\PassOptionsToPackage{sols}{handout}
\documentclass[main]{subfiles}
\setcounterrefbeforedot{chapter}{m-ws5-3}
\setcounterrefafterdot{section}{m-ws5-3}
\addtocounter{section}{-1}
\usepackage{solutions}

\begin{document}

\ifSubfilesClassLoaded{%
  \chaptermark{\chapterfive}
}{}

\section{Definition of the Definite Integral} \label{ws5-3}

\begin{problemonly}
    \begin{framed}

\textbf{Key Ideas}:

\begin{itemize}[topsep=0pt,partopsep=0pt,parsep=8pt,itemsep=0pt]

\item You should be able to write a general Riemann sum (that is, a
  left-hand sum or right-hand sum with $n$ rectangles) using the notation
  $\Delta x$ and $x_0, \ldots, x_n$.  You should be able to explain what
  $\Delta x$ and $x_0, \ldots, x_n$ are through a picture and by giving
  formulas for them.

\item You should understand how we use $\sum$ notation to express Riemann
  sums.  You should be aware that the  ``index'' variable in $\sum$
  notation is a dummy variable, so $\displaystyle \sum_{k=1}^n f(x_k)
  \Delta x$, $\displaystyle \sum_{i=1}^n f(x_i) \Delta x$, and
  $\displaystyle \sum_{j=1}^n f(x_j) \Delta x$ all mean
  the same thing.

\item You should understand the definition of the definite integral as a
  limit of Riemann sums and be able to apply this definition to compute a
  definite integral.  (You should also see why this is not a preferred way
  of evaluating an integral!)

\item You should develop flexibility with our 3 ways of thinking about the
  definite integral (its definition as a limit of Riemann sums, its
  graphical interpretation as a signed area, and its interpretation as net
  change if the integrand is a rate function). 

\end{itemize}
\end{framed}
\end{problemonly}

\begin{enumerate}
\item \emph{Recap!}
  \begin{enumerate}
  \item
    \begin{problem}
      We have 3 main ways of thinking about the definite integral
      $\displaystyle \int_a^b f(x)\,dx$.  What are they?
    \end{problem}

      \begin{itemize}
      \item
        \solonly{
        If $a < b$, we visualize $\displaystyle \int_a^b f(x)\,dx$
        as the signed area under $f$ from $a$ to $b$.
        }
        \pspace{36pt}
      \item
        \solonly{
        If $f$ is a rate function, then we interpret $\displaystyle
        \int_a^b f(x)\,dx$ as the net change in amount from $x = a$ to
        $x = b$.  For instance, if $f(x)$ represents the rate at which
        water is entering or leaving a pool at time $x$, then
        $\displaystyle \int_a^b f(x)\,dx$ is the net change in amount of
        water from time $a$ to time $b$, or (amount of water at time
        $b$) $-$ (amount of water at time $a$).
        }
        \pspace{36pt}
      \item
        \solonly{
        $\displaystyle \int_a^b f(x)\,dx$ is defined to be
        $\displaystyle \lim_{n \to \infty} R_n$ or $\displaystyle
        \lim_{n \to \infty} L_n$.
        }
        \pspace{36pt}
      \end{itemize}
        \pspace{36pt}
  \item
    \begin{problem}[\vfill]
      So far, when are we able to evaluate $\displaystyle \int_a^b
      f(x)\,dx$ exactly?
    \end{problem}

    \begin{solution}
      So far, we've only been able to do this when the graph of $f$ on
      $[a, b]$ has been simple geometrically (made up of quarter-circles
      or line segments) or if it has some particularly nice symmetry that
      allows us to conclude the integral is 0.
    \end{solution}

  \end{enumerate}

\item 

  \note{The point of this problem is for students to figure out how to
    calculate $\Delta x$ and $x_0, x_1, \ldots, x_n$ in a very concrete
    example before we generalize.}

  \begin{problem}[0.25in]
    Sketch the right-hand sum $R_4$ that
    approximates $\displaystyle \int_4^7 f(x)\,dx$, and write out $R_4$ as
    a sum.

    \begin{tikzpicture}[declare function={f(\x)=-1685 + 13268*\x -
        3456*\x^2 + 256*\x^3;}] 
      \begin{axis}[ymin=0, hide y axis, xtick={4, 7}, width=3.5in]
        \node[yshift=20pt] at (6,{f(6)}) {$y=f(x)$};
        \addplot+[domain=3:8] {f(x)};
        % \draw[dashed] (4,{f(4)}) -- (4,0);
        % \draw[dashed] (7,{f(7)}) -- (7,0);
        \solonly{
        \addplot+[ybar interval, plusarea] coordinates {
            (4, {f(4.75)}) (4.75, {f(5.5)}) (5.5, {f(6.25)}) (6.25, {f(7)})
            (7, {f(7)})
          }; %
        }
      \end{axis}
    \end{tikzpicture}
  \end{problem}

  \begin{solution}
    The interval $[4, 7]$ is 3 units wide, so if we're going to use
    4 rectangles, each rectangle should be $\frac{3}{4} = 0.75$ units
    wide.
    Therefore, $R_4 = \boxed{0.75 f(4.75) + 0.75 f(5.5) + 0.75 f(6.25)+
    0.75 f(7)}$. 
  \end{solution}

\item 
  {
    \providecommand{\graphcommand}{}
    \renewcommand{\graphcommand}[2]{%
      \begin{tikzpicture}[declare function={f(\x)=-16.85 + 132.68*\x -
          34.56*\x^2 + 2.56*\x^3;}] 
        \begin{axis}[ymin=0, hide y axis, xtick={4, 7}, xticklabels={$a$,
            $b$}, width=2.25in, cycle list=black, #1]
          \addplot+[domain=4:7, fill=cyan, draw=none] {f(x)} \closedcycle;
          #2
          \addplot+[domain=3:8] {f(x)};
          \addplot+[domain=3:8, thin] {0};
          \node[yshift=15pt] at (6,{f(6)}) {$y=f(x)$};
        \end{axis}
      \end{tikzpicture} 
    }

    Suppose we have a function $y = f(x)$ and want to estimate
    $\displaystyle \int_a^b f(x)\,dx$ using $R_n$, the right-hand sum with
    $n$ rectangles.

    \begin{enumerate}
    \item

      \note{The idea of $x_k$ and the ``$k$-th rectangle'' can be confusing
        to students (they wonder, ``what's $k$?'').}

      We start by \textbf{slicing} the area into $n$ pieces of equal width
      $\Delta x$ (pictured in Figure 2 with $n = 6$ and in Figure 3 with $n
      = 12$).
      \begin{center}
        \begin{tabular}{c @{\hspace{0.35in}} c @{\hspace{0.35in}} c}
          \graphcommand{}{} &
          \graphcommand{xtick={4, 4.5, ..., 7}, xticklabels={$x_0$, $x_1$,
              $x_2$, $x_3$, $x_4$, $x_5$, $x_6$}, tick label
            style={font=\scriptsize}}{ 
            \pgfplotsforeachungrouped \x in {4, 4.5, ..., 7}
            {
              \pgfmathsetmacro{\y}{f(\x)}
              \edef\temp{%
                \noexpand\draw[red] (axis cs: \x, 0) -- (axis cs: \x, \y);
              }
              \temp
            } 
          } &
          \graphcommand{xtick={4, 7}, xticklabels={$x_0$, $x_n$}, tick label
            style={font=\scriptsize}}{ 
            \pgfplotsforeachungrouped \x in {4, 4.25, ..., 7}
            {
              \pgfmathsetmacro{\y}{f(\x)}
              \edef\temp{%
                \noexpand\draw[red] (axis cs: \x, 0) -- (axis cs: \x, \y);
              }
              \temp
            } 
          }      
          \\    
          Figure 1 & Figure 2 & Figure 3 
        \end{tabular}
      \end{center}
        \pspace{8pt}
      \begin{problem}
        What is $\Delta x$ in terms of $a$, $b$, and $n$?
        $\fitb{1in}{\frac{b-a}{n}}$ \pspace{8pt}
      \end{problem}

      \begin{problem}
        You can see that we label the values where we slice as $x_0, x_1,
        \ldots, x_n$. \pspace{8pt}
        
        What is $x_0$ (in terms of $a, b, \Delta x$)?
        $\fitb{0.5in}{a}$ \pspace{8pt}

        What is $x_1$? $\fitb{1in}{a + \Delta x}$ $x_2$?  $\fitb{1in}{a +
          2 \Delta x}$ $x_5$? $\fitb{1in}{a + 5 \Delta x}$ \pspace{8pt}

        In general, what is $x_k$? $\fitb{1in}{a + k \Delta x}$
      \end{problem}

    \pspace{\newpage}
    
    \item Since we are doing a right-hand sum, we \textbf{approximate} the
      area of each rectangle using a rectangle whose height is the height of
      the right side of the slice, as shown.

      \graphcommand{xtick={4, 4.5, ..., 7}, xticklabels={$x_0$, $x_1$,
          $x_2$, $x_3$, $x_4$, $x_5$, $x_6$}, tick label
        style={font=\scriptsize}}{
        \addplot+[ybar interval, fill=purple, draw=yellow] coordinates
        {
          (4, {f(4.5)}) (4.5, {f(5)}) (5, {f(5.5)}) (5.5, {f(6)}) (6,
          {f(6.5)}) (6.5, {f(7)}) (7, {f(7)})
        };
      }
      \hspace{0.25in}
      \graphcommand{xtick={4, 7}, xticklabels={$x_0$, $x_n$}, tick label
        style={font=\scriptsize}}{ 
        \addplot+[ybar interval, fill=purple, draw=yellow] coordinates
        {
          (4, {f(4.25)}) (4.25, {f(4.5)}) (4.5, {f(4.75)}) (4.75, {f(5)})
          (5, {f(5.25)}) (5.25, {f(5.5)}) (5.5, {f(5.75)}) (5.75, {f(6)})
          (6, {f(6.25)}) (6.25, {f(6.5)}) (6.5, {f(6.75)}) (6.75, {f(7)})
          (7, {f(7)})  
        };
      }      
        \pspace{16pt}
        
      \begin{problem}
        What is the signed area of the 1st rectangle? $\fitb{1.5in}{f(x_1)
          \Delta x}$ 
      \end{problem}
      \pspace{16pt}

      \begin{problem}
        What is the signed area of the 4th rectangle? $\fitb{1.5in}{f(x_4)
          \Delta x}$ 
      \end{problem}
      \pspace{16pt}

      \begin{problem}
        What is the signed area of the $k$-th rectangle?
        $\fitb{1.5in}{f(x_k) \Delta x}$  
      \end{problem}
      \pspace{16pt}
    
    \item

      \note{I'll talk here about using sigma notation to express the sum.  I
        like to just have them think about it as, ``add up the signed area
        of the $k$-th rectangle'' where $k$ goes from $1$ to $n$ to get the
        sum.} 

      \begin{problem}[36pt]
        What is $R_n$?
      \end{problem}

      \begin{solution}
        $f(x_1) \Delta x+ f(x_2) \Delta x+ f(x_3) \Delta x+ ... +
        f(x_{n-1}) \Delta x+ f(x_n) \Delta x$, which can be written in
        $\sum$ notation as $\displaystyle \sum_{k=1}^n f(x_{k}) \Delta x$
      \end{solution}

    \item

      \note{You can write $L_n$ in sigma notation as either $\displaystyle
        \sum_{k=1}^n f(x_{k-1}) \Delta x$ or $\displaystyle
        \sum_{k=0}^{n-1} f(x_k) \Delta x$; it doesn't matter.}

      \begin{problem}[\vfill]
        What if you instead wanted to use a left-hand sum?  What would the
        signed area of the $k$-th rectangle be?  What is $L_n$?
      \end{problem}

      \begin{solution}
        The left-hand sum is $L_n = f(x_0) \Delta x + f(x_1) \Delta x +
        \cdots + f(x_{n-1}) \Delta x$, which can be written in $\sum$
        notation as either $\displaystyle \sum_{k=0}^{n-1} f(x_k) \Delta
        x$ or $\displaystyle \sum_{k=1}^n f(x_{k-1}) \Delta x$. 
      \end{solution}

    \end{enumerate}
  }

\end{enumerate}
\begin{problemonly}
  \begin{framed}
    \textbf{The official definition of the definite integral.}  The
    \uline{definite integral of $f(x)$ from $x = a$ to $x = b$}, denoted
    $\displaystyle \int_a^b f(x)\,dx$, is defined to be
    \[ \lim_{n \to \infty} \underbrace{\sum_{k=1}^n f(x_k)\, \Delta x}_{R_n}
    = \lim_{n \to \infty} \underbrace{\sum_{k=0}^{n-1} f(x_k)\, \Delta
      x}_{L_n}, \text{ where } \Delta x = \frac{b - a}{n} \text{ and } x_k
    = a + k \Delta x\] provided these limits exist.

    \textit{If $f$ is continuous on $[a, b]$, then the limits always
      exist.}
  \end{framed}
\end{problemonly}

\pspace{\newpage}

\begin{enumerate}[resume]
\item 
  \note{Calculating $R_n$ in this example is relatively simple
    algebraically (and therefore is not representative of the work they'd
    need to do for most integrals).  That's okay; we don't expect them to
    become fluent with using the definition to evaluate integrals; we just
    want to give them a sense of what's involved and why we desperately
    want another way to evaluate integrals!}

  In this problem, we'll evaluate $\displaystyle \int_0^1 x^2\,dx$ using
  the definition of the definite integral as $\displaystyle \lim_{n \to
    \infty} R_n$. 

  \begin{enumerate}
  \item
    \note{This is very challenging, so we'll definitely do it together.}

    \begin{problem}[\vfill]
      Write a formula for $R_n$ and simplify it.  The following fact will
      be useful.
      \[ 1^2 + 2^2 + 3^2 + \cdots + n^2, \text{ which can also be written
        as } \sum_{k=1}^n k^2, \text{ is equal to } \frac{2n^3 + 3n^2 +
        n}{6}\]
    \end{problem}

    \begin{solution}
      First, the interval $[0, 1]$ has width $1$, so if we're using $n$
      rectangles, then each rectangle should have width $\Delta x =
      \frac{1}{n}$.  When we slice the interval $[0, 1]$ into $n$ pieces
      of equal width, the $x$-values we slice at are $x_0 = 0$, $x_1 =
      \frac{1}{n}$, $x_2 = \frac{2}{n}, \ldots, x_n = \frac{n}{n} = 1$:
      \begin{center}
        \begin{tikzpicture}[declare function={f(\x) = \x^2;}]
          \begin{axis}[cycle list=black, xtick={0, 0.125, ..., 1.01},
            xticklabels={$x_0$, $x_1$, $x_2$, $x_3$, $\cdots$ , , , , $x_n$}, 
            ytick=\empty, axis x line=bottom, ymin=0, clip=false] 
            \addplot+[ybar interval, plusarea] coordinates { (0, {f(1/8)})
              (1/8, {f(2/8)}) (2/8, {f(3/8)}) (3/8, {f(4/8)}) (4/8,
              {f(5/8)}) (5/8, {f(6/8)}) (6/8, {f(7/8)}) (7/8, {f(1)}) (1,
              {f(1)}) }; %
            \addplot+[domain=0:1.1] {f(x)}; %
            \pgfplotsforeachungrouped \x in {0, 0.125, 0.25, 0.375, 1} %
            {%
              \edef\temp{%
                \noexpand\draw (axis cs: \x, 0) node[yshift=-13pt, rotate=90]
                {$=$};                
              }%
              \temp
            }
            \draw (axis cs: 0, 0) node[font=\footnotesize, yshift=-21pt] {$0$};
            \draw (axis cs: 0.125, 0) node[font=\footnotesize, yshift=-24pt] {$\frac{1}{n}$};
            \draw (axis cs: 0.25, 0) node[font=\footnotesize, yshift=-24pt] {$\frac{2}{n}$};
            \draw (axis cs: 0.375, 0) node[font=\footnotesize, yshift=-24pt] {$\frac{3}{n}$};
            \draw (axis cs: 1, 0) node[font=\footnotesize, yshift=-24pt] {$\frac{n}{n}$};
          \end{axis}
        \end{tikzpicture}
      \end{center}
      If we write $f(x) = x^2$, then 
      \begin{align*}
        R_n &= f(x_1) \Delta x + f(x_2) \Delta x + f(x_3) \Delta x +
        \cdots + f (x_n) \Delta x \\
        &= f \left( \frac{1}{n} \right) \frac{1}{n} + f \left( \frac{2}{n}
        \right) \frac{1}{n} + f \left( \frac{3}{n} \right) \frac{1}{n} +
        \cdots + f \left( \frac{n}{n} \right) \frac{1}{n} \\
        &= \left( \frac{1}{n} \right)^2 \frac{1}{n} + \left( \frac{2}{n}
        \right)^2 \frac{1}{n} + \left( \frac{3}{n} \right)^2 \frac{1}{n} +
        \cdots + \left( \frac{n}{n} \right)^2 \frac{1}{n} \\
        &= \frac{1}{n^3} + \frac{2^2}{n^3} + \frac{3^2}{n^3} + \cdots +
        \frac{n^2}{n^3} \\
        &= \frac{1}{n^3} \left( 1 + 2^2 + 3^2 + \cdots + n^2 \right) \\
        \intertext{Using the useful fact given in the problem, this is
          equal to}
        &= \frac{1}{n^3} \cdot \frac{2n^3 + 3n^2 + n}{6} \\
        &= \boxed{\frac{2n^3 + 3n^2 + n}{6n^3}}        
      \end{align*}
    \end{solution}

  % \item

  %   \note{Here, I just want to make sure they understand what they figured
  %     out in the previous part.  (In other words, I want them to realize
  %     that they can just plug $n = 4$ into their answer to
  %     {{{{\#4}}(a)}}.)}

  %   \begin{problem}[0.5in]
  %     What is $R_4$?  (Simplify your answer.)
  %   \end{problem}

  %   \begin{solution}
  %     In {{{{\#4}}(a)}}, we found that
  %     $\displaystyle R_n = \frac{2n^3 + 3n^2 + n}{6n^3}$; plugging in $n =
  %     4$ gives $\displaystyle R_4 = \frac{2 \cdot 4^3 + 3 \cdot 4^2 + 4}{6
  %       \cdot 4^3} = \boxed{\frac{15}{32}}$. 
  %   \end{solution}

  \item    
    \begin{problem}[36pt]
      What is $\displaystyle \int_0^1 x^2\,dx$?
    \end{problem}

    \begin{solution}
      In {{{{\#4}}(a)}}, we found that
      $\displaystyle R_n = \frac{2n^3 + 3n^2 + n}{6n^3}$.  By the
      definition of the definite integral, $\displaystyle \int_0^1 x^2\,dx
      = \lim_{n \to \infty} R_n = \lim_{n \to \infty} \frac{2n^3 + 3n^2 +
        n}{6n^3}$.  This is a $\frac{\infty}{\infty}$ limit; in the
      numerator, the fastest growing term is $2n^3$, so this limit is
      equal to $\displaystyle \lim_{n \to \infty} \frac{2n^3}{6n^3} =
      \boxed{\frac{1}{3}}$.
    \end{solution}

  \end{enumerate}
    \pspace{\newpage}
    
\item {
  \note{This problem encourages students to generalize their thinking
    about left-hand and right-hand sums a bit.}

  Suppose we want to approximate $\displaystyle
  \int_0^{1.5} \sin x \,dx$ using various left-hand and right-hand sums.

  \begin{enumerate}
  \item 
      Put the following expressions in ascending order.
      \[ 
      \int_0^{1.5} \sin x \,dx \hspace{0.5in} L_4 \hspace{0.5in} R_4
      \hspace{0.5in} L_{20} \hspace{0.5in} R_{20} \hspace{0.5in} L_{100}
      \hspace{0.5in} R_{100}
      \]
      \pspace{16pt}
      Which of the three diagrams below shows $L_4$? Which one
      shows $R_4$? Which one shows $|R_4 - L_4|$?
    
    \newcommand{\graphcommand}[2]{
    \begin{tikzpicture}[declare function={f(\x)=sin(deg(\x));}]
        \begin{axis}[
        xmax=1.75, 
        ytick=\empty,
        xtick={0,1.5},
        minor tick num = 3,
        minor tick length = 5pt,
        width = 2.2in,
        clip = false,
        cycle list=black,
        #2]
            \addplot[domain=0:1.5] {f(x)};
            #1
        \end{axis}
    \end{tikzpicture}
    }
    \begin{tabular}{ccc}
        \graphcommand{
        \addplot+[ybar interval, pattern=flexible hatch, pattern color=blue,
        hatch direction=-1] coordinates {
        (0,{f(0.375)}) (0.375,{f(0.75)}) (0.75,{f(1.125)}) (1.125,{f(1.5)})
        (1.5,{f(1.5)})
        };
        }{\solonly{title=$R_4$}}
        &
        \graphcommand{
        \addplot+[ybar interval, pattern=flexible hatch, pattern color=red]
        coordinates {(0,{f(0)}) 
        (0.375,{f(0.375)}) (0.75,{f(0.75)}) (1.125,{f(1.125)})
        (1.5,{f(1.125)})};
        }{\solonly{title=$L_4$}}
        &
        % \graphcommand{
        % \filldraw[fill=green, fill opacity=0.5, draw=green!70!black] 
        % (0,{f(0)}) rectangle (0.375,{f(0.375)})
        % (0.375,{f(0.375)}) rectangle (0.75,{f(0.75)})
        % (0.75,{f(0.75)}) rectangle (1.125,{f(1.125)})
        % (1.125,{f(1.125)}) rectangle (1.5,{f(1.5)});
        % }{}
        \graphcommand{
        \addplot+[ybar interval, pattern=flexible hatch, pattern color=red]
        coordinates {(0,{f(0)}) 
        (0.375,{f(0.375)}) (0.75,{f(0.75)}) (1.125,{f(1.125)})
        (1.5,{f(1.125)})};
        \addplot+[ybar interval, pattern=flexible hatch, pattern color=blue,
        hatch direction=-1] coordinates {
        (0,{f(0.375)}) (0.375,{f(0.75)}) (0.75,{f(1.125)}) (1.125,{f(1.5)})
        (1.5,{f(1.5)})
        };
        }{\solonly{title=$|R_4-L_4|$}} 
    \end{tabular}

    \begin{solution}
      (This problem is very similar to
          {{{{\#4}} on {the worksheet ``The Definite Integral''}}}.) 
          Since $f(x) = \sin x$
      is increasing on $[0, 1.5]$, left-hand sums will be underestimates
      for $\displaystyle \int_0^{1.5} f(x)\,dx$, and right-hand sums will
      be overrestimates.

      Here are pictures of $L_4$ and $L_{20}$:
      \begin{center}
        \begin{tikzpicture}[declare function={f(\x)=sin(deg(\x));}]
          \begin{axis}[ytick=\empty, xtick={1.5}, cycle list=black,
            clip=false] 
            \addplot+[ybar interval, plusarea] coordinates {
              (0, {f(0)}) (3/8, {f(3/8)}) (3/4, {f(3/4)}) (9/8, {f(9/8)})
              (3/2, {f(9/8)}) };
            \addplot+[domain=0:1.7] {f(x)};
            \draw (axis cs: -0.5, 0.5) node {$L_4$};
          \end{axis}
        \end{tikzpicture}

        \begin{tikzpicture}[declare function={f(\x)=sin(deg(\x));}]
          \begin{axis}[ytick=\empty, xtick={1.5}, cycle list=black,
            clip=false] 
            \addplot+[ybar interval, plusarea] coordinates {
              (0, {f(0)}) (3/40, {f(3/40)}) (3/20, {f(3/20)}) (9/40, {f(9/40)}) (3/10, {f(3/10)}) (3/8, {f(3/8)}) (9/20, {f(9/20)}) (21/40, {f(21/40)}) (3/5, {f(3/5)}) (27/40, {f(27/40)}) (3/4, {f(3/4)}) (33/40, {f(33/40)}) (9/10, {f(9/10)}) (39/40, {f(39/40)}) (21/20, {f(21/20)}) (9/8, {f(9/8)}) (6/5, {f(6/5)}) (51/40, {f(51/40)}) (27/20, {f(27/20)}) (57/40, {f(57/40)}) (3/2, {f(57/40)})  };
            \addplot+[domain=0:1.7] {f(x)};
            \draw (axis cs: -0.5, 0.5) node {$L_{20}$};
          \end{axis}
        \end{tikzpicture}
      \end{center}
      As we can see, $L_{20}$ is much closer to $\displaystyle
      \int_0^{1.5} f(x)\,dx$ than $L_4$ is.  By the same idea, $L_{100}$
      will be even closer to $\displaystyle \int_0^{1.5} f(x)\,dx$ than
      $L_{20}$ is.  Similarly, among the right-hand sums, $R_{100}$ is
      closest to $\displaystyle \int_0^{1.5} f(x)\,dx$, followed by
      $R_{20}$ and then finally by $R_4$.  That gives us the following
      ordering:
      \[ \boxed{L_4 < L_{20} < L_{100} < \int_0^{1.5} f(x)\,dx < R_{100} <
        R_{20} < R_4}.\]
    \end{solution}

  \item
    \begin{problem}[\vfill]
      Find $|R_4 - L_4|$.  (What's the least
      work you can do to find $|R_4 - L_4|$?  Do you need to calculate
      $R_4$ and $L_4$?)
    \end{problem}

    \begin{solution}
      Here are pictures of $L_4$ (in red stripes) and $R_4$ (in blue
      stripes): 
      \begin{center}
        \begin{tikzpicture}[declare function={f(\x)=sin(deg(\x));}]
          \begin{axis}[ytick=\empty, xtick={1.5}, cycle list=black,
            clip=false, xtick={0, 0.375, ..., 1.51}, xticklabels={$x_0$,
              $x_1$, $x_2$, $x_3$, $x_4$}] 
            \addplot+[ybar interval, pattern=flexible hatch, pattern color=red] coordinates {
              (0, {f(0)}) (3/8, {f(3/8)}) (3/4, {f(3/4)}) (9/8, {f(9/8)})
              (3/2, {f(9/8)}) };
            \addplot+[ybar interval, pattern=flexible hatch, pattern color=blue, hatch direction=-1] coordinates {
              (0, {f(3/8)}) (3/8, {f(3/4)}) (3/4, {f(9/8)}) (9/8,
              {f(3/2)}) (3/2, {f(3/2)})  }; 
            \addplot+[domain=0:1.7] {f(x)};
            \draw (axis cs: 0, 0) node[font=\footnotesize, fill=white,
            inner sep=1pt, below, yshift=-2pt] {$x_0$};
            \draw (axis cs: 1.5, 0) node[font=\footnotesize, yshift=-12pt, rotate=90] {$=$};
            \draw (axis cs: 1.5, 0) node[font=\footnotesize, yshift=-20pt] {$1.5$};
          \end{axis}
        \end{tikzpicture}
      \end{center}
      Notice that the first 3 rectangles in $R_4$ are the same as the last
      3 rectangles in $L_4$; that is,
      \[ R_4 - \text{(last
        rectangle in $R_4$)} = L_4 - \text{(first rectangle in $L_4$)}.\]
      Rearranging this equation,
      \begin{align*}
        R_4 - L_4 &= \text{(last rectangle in $R_4$)} - \text{(first
          rectangle in $L_4$)} \\
        &= f(1.5) \Delta x - f(0) \Delta x \text{ where $\Delta x =
          \frac{1.5}{4}$ is the width of each rectangle} \\
        \shortintertext{So,}
        |R_4 - L_4| &= |f(1.5) - f(0)| \Delta x \\
        &= \left| \sin 1.5 - \sin 0 \right| \frac{1.5}{4} \\
        &= \boxed{\frac{1.5 \sin 1.5}{4}}
      \end{align*}
      \emph{Alternate Method:} If you prefer, you could write out both
      $L_4$ and $R_4$:
      \begin{align*}
        R_4 &= f(x_1) \Delta x + f(x_2) \Delta x + f(x_3) \Delta x +
        f(x_4) \Delta x \\
        L_4 &= f(x_0) \Delta x + f(x_1) \Delta x + f(x_2) \Delta x +
        f(x_3) \Delta x 
      \end{align*}
      Subtracting, $R_4 - L_4 = f(x_4) \Delta x - f(x_0) \Delta x = f(1.5)
      \Delta x - f(0) \Delta x$, as we said before.
    \end{solution}

  \item
    \begin{problem}[\stretch{0.5}]
      Find $|R_{100} - L_{100}|$.    
    \end{problem}    

    \begin{solution}
      By the same reasoning as in
      {{{{\#5}}(b)}}, $|R_{100} - L_{100}| =
      |(\sin 1.5) \Delta x - (\sin 0) \Delta x|$ where now $\Delta x =
      \frac{1.5}{100}$ (it's the width of each rectangle if we use 100
      rectangles).  This can be rewritten as $\boxed{ \frac{1.5 \sin
          1.5}{100}}$.
    \end{solution}

  \item

    \begin{problem}[\vfill]
      How large do we need $n$ to be ensure $|R_n - L_n| < 0.05$?  Why
      might we care to have $|R_n - L_n| < 0.05?$
    \end{problem}

    \begin{solution}
      Using the same reasoning as in
      {{{{\#5}}(b)}} and
      {{{{\#5}}(c)}}, $|R_n - L_n| =
      \frac{1.5 \sin 1.5}{n}$, so we want 
      $\frac{1.5 \sin 1.5}{n} < 0.05$.  Rearranging this, we need $n >
      \frac{1.5 \sin 1.5}{0.05} = 30 \sin 1.5$.

      Since $1.5$ is just a little bit less than $\frac{\pi}{2}$, $\sin
      1.5$ is just a bit less than $1$.  Therefore, if we take $n$ to be
      \fbox{at least $30$}, we can be sure that $n > 30 \sin 1.5$. 

      As we said in {{{{\#5}}(a)}}, left-hand sums
      are underestimates for $\displaystyle \int_0^{1.5} \sin x\,dx$ and
      right-hand sums are overestimates, so we always have $\displaystyle
      L_n < \int_0^{1.5} \sin x\,dx < R_n$.  If $R_n$ and $L_n$ are within
      $0.05$ of each other, that means we have $\displaystyle \int_0^{1.5}
      \sin x\,dx$ pinned between two values that are within $0.05$ of each
      other.  This means that we can estimate $\displaystyle \int_0^{1.5}
      \sin x\,dx$ and be certain that we are off by less than $0.05$.

      (For example, $L_{30} = 0.904132$ and $R_{30} = 0.954007$, so we can
      be sure that $\displaystyle \int_0^{1.5} \sin x\,dx$ is between
      these two values.)
    \end{solution}

  \end{enumerate}
}
\end{enumerate}

\end{document}